\section{待ち行列}

\subsection{令和元年度 問3-4}
\label{q:r01-reg-basic-3-4}

\noindent\textbf{原典：} sources/R01/basic.pdf、PDF物理13頁。

\textbf{公式正答：} 4（single）。

\par\textit{原典本文の校正と独自解説は未作成。分類のみの登録項目。}

\section{ベクトル・微分演算}

\subsection{平成23年度 問3-5}
\label{q:h23-reg-basic-3-5}

\noindent\textbf{原典：} sources/H23/basic.pdf、PDF物理11頁。

\textbf{公式正答：} 1（single）。

\par\textit{原典本文の校正と独自解説は未作成。分類のみの登録項目。}

\subsection{平成25年度 問3-6}
\label{q:h25-reg-basic-3-6}

\noindent\textbf{原典：} sources/H25/basic.pdf、PDF物理14頁。

\textbf{公式正答：} 4（single）。

\par\textit{原典本文の校正と独自解説は未作成。分類のみの登録項目。}

\subsection{平成26年度 問3-1}
\label{q:h26-reg-basic-3-1}

\noindent\textbf{原典：} sources/H26/basic.pdf、PDF物理9頁。

\textbf{公式正答：} 2（single）。

\par\textit{原典本文の校正と独自解説は未作成。分類のみの登録項目。}

\subsection{平成26年度 問3-4}
\label{q:h26-reg-basic-3-4}

\noindent\textbf{原典：} sources/H26/basic.pdf、PDF物理10頁。

\textbf{公式正答：} 1（single）。

\par\textit{原典本文の校正と独自解説は未作成。分類のみの登録項目。}

\subsection{平成28年度 問3-3}
\label{q:h28-reg-basic-3-3}

\noindent\textbf{原典：} sources/H28/basic.pdf、PDF物理13頁。

\textbf{公式正答：} 2（single）。

\par\textit{原典本文の校正と独自解説は未作成。分類のみの登録項目。}

\subsection{平成29年度 問3-2}
\label{q:h29-reg-basic-3-2}

\noindent\textbf{原典：} sources/H29/basic.pdf、PDF物理11頁。

\textbf{公式正答：} 4（single）。

\par\textit{原典本文の校正と独自解説は未作成。分類のみの登録項目。}

\subsection{平成30年度 問3-2}
\label{q:h30-reg-basic-3-2}

\noindent\textbf{原典：} sources/H30/basic.pdf、PDF物理11頁。

\textbf{公式正答：} 4（single）。

\par\textit{原典本文の校正と独自解説は未作成。分類のみの登録項目。}

\subsection{令和元年度 問3-1}
\label{q:r01-reg-basic-3-1}

\noindent\textbf{原典：} sources/R01/basic.pdf、PDF物理12頁。

\textbf{公式正答：} 5（single）。

\par\textit{原典本文の校正と独自解説は未作成。分類のみの登録項目。}

\subsection{令和2年度 問3-1}
\label{q:r02-reg-basic-3-1}

\noindent\textbf{原典：} sources/R02/basic.pdf、PDF物理11頁。

\textbf{公式正答：} 2（single）。

\par\textit{原典本文の校正と独自解説は未作成。分類のみの登録項目。}

\subsection{令和2年度 問3-2}
\label{q:r02-reg-basic-3-2}

\noindent\textbf{原典：} sources/R02/basic.pdf、PDF物理11頁。

\textbf{公式正答：} 4（single）。

\par\textit{原典本文の校正と独自解説は未作成。分類のみの登録項目。}

\subsection{令和3年度 問3-1}
\label{q:r03-reg-basic-3-1}

\noindent\textbf{原典：} sources/R03/basic.pdf、PDF物理11頁。

\textbf{公式正答：} 4（single）。

\par\textit{原典本文の校正と独自解説は未作成。分類のみの登録項目。}

\subsection{令和4年度 問3-2}
\label{q:r04-reg-basic-3-2}

\noindent\textbf{原典：} sources/R04/basic.pdf、PDF物理10頁。

\textbf{公式正答：} 2（single）。

\par\textit{原典本文の校正と独自解説は未作成。分類のみの登録項目。}

\subsection{令和6年度 問3-2}
\label{q:r06-reg-basic-3-2}

\noindent\textbf{原典：} sources/R06/basic.pdf、PDF物理10頁。

\textbf{公式正答：} 2（single）。

\par\textit{原典本文の校正と独自解説は未作成。分類のみの登録項目。}

\subsection{令和7年度 問3-2}
\label{q:r07-reg-basic-3-2}

\noindent\textbf{原典：} sources/R07/basic.pdf、PDF物理8頁。

\textbf{公式正答：} 1（single）。

\par\textit{原典本文の校正と独自解説は未作成。分類のみの登録項目。}

\section{微積分・求積}

\subsection{平成24年度 問3-3}
\label{q:h24-reg-basic-3-3}

\noindent\textbf{原典：} sources/H24/basic.pdf、PDF物理9頁。

\textbf{公式正答：} 1（single）。

\par\textit{原典本文の校正と独自解説は未作成。分類のみの登録項目。}

\subsection{平成24年度 問3-4}
\label{q:h24-reg-basic-3-4}

\noindent\textbf{原典：} sources/H24/basic.pdf、PDF物理10頁。

\textbf{公式正答：} 4（single）。

\par\textit{原典本文の校正と独自解説は未作成。分類のみの登録項目。}

\subsection{平成26年度 問3-3}
\label{q:h26-reg-basic-3-3}

\noindent\textbf{原典：} sources/H26/basic.pdf、PDF物理10頁。

\textbf{公式正答：} 5（single）。

\par\textit{原典本文の校正と独自解説は未作成。分類のみの登録項目。}

\subsection{平成28年度 問3-1}
\label{q:h28-reg-basic-3-1}

\noindent\textbf{原典：} sources/H28/basic.pdf、PDF物理12頁。

\textbf{公式正答：} 3（single）。

\par\textit{原典本文の校正と独自解説は未作成。分類のみの登録項目。}

\subsection{平成28年度 問3-2}
\label{q:h28-reg-basic-3-2}

\noindent\textbf{原典：} sources/H28/basic.pdf、PDF物理12頁。

\textbf{公式正答：} 3（single）。

\par\textit{原典本文の校正と独自解説は未作成。分類のみの登録項目。}

\subsection{平成28年度 問3-4}
\label{q:h28-reg-basic-3-4}

\noindent\textbf{原典：} sources/H28/basic.pdf、PDF物理14頁。

\textbf{公式正答：} 5（single）。

\par\textit{原典本文の校正と独自解説は未作成。分類のみの登録項目。}

\subsection{平成29年度 問3-1}
\label{q:h29-reg-basic-3-1}

\noindent\textbf{原典：} sources/H29/basic.pdf、PDF物理11頁。

\textbf{公式正答：} 5（single）。

\par\textit{原典本文の校正と独自解説は未作成。分類のみの登録項目。}

\subsection{令和元年度 問3-2}
\label{q:r01-reg-basic-3-2}

\noindent\textbf{原典：} sources/R01/basic.pdf、PDF物理12頁。

\textbf{公式正答：} 5（single）。

\par\textit{原典本文の校正と独自解説は未作成。分類のみの登録項目。}

\subsection{令和元年度再試験 問3-1}
\label{q:r01-ret-basic-3-1}

\noindent\textbf{原典：} sources/R01\_retest/basic.pdf、PDF物理11頁。

\textbf{公式正答：} 5（single）。

\par\textit{原典本文の校正と独自解説は未作成。分類のみの登録項目。}

\subsection{令和元年度再試験 問3-4}
\label{q:r01-ret-basic-3-4}

\noindent\textbf{原典：} sources/R01\_retest/basic.pdf、PDF物理12頁。

\textbf{公式正答：} 5（single）。

\par\textit{原典本文の校正と独自解説は未作成。分類のみの登録項目。}

\subsection{令和3年度 問3-2}
\label{q:r03-reg-basic-3-2}

\noindent\textbf{原典：} sources/R03/basic.pdf、PDF物理11頁。

\textbf{公式正答：} 2（single）。

\par\textit{原典本文の校正と独自解説は未作成。分類のみの登録項目。}

\subsection{令和4年度 問3-1}
\label{q:r04-reg-basic-3-1}

\noindent\textbf{原典：} sources/R04/basic.pdf、PDF物理10頁。

\textbf{公式正答：} 4（single）。

\par\textit{原典本文の校正と独自解説は未作成。分類のみの登録項目。}

\subsection{令和5年度 問3-2}
\label{q:r05-reg-basic-3-2}

\noindent\textbf{原典：} sources/R05/basic.pdf、PDF物理9頁。

\textbf{公式正答：} 2（single）。

\par\textit{原典本文の校正と独自解説は未作成。分類のみの登録項目。}

\subsection{令和6年度 問3-1}
\label{q:r06-reg-basic-3-1}

\noindent\textbf{原典：} sources/R06/basic.pdf、PDF物理10頁。

\textbf{公式正答：} 4（single）。

\par\textit{原典本文の校正と独自解説は未作成。分類のみの登録項目。}

\subsection{令和7年度 問3-3}
\label{q:r07-reg-basic-3-3}

\noindent\textbf{原典：} sources/R07/basic.pdf、PDF物理9頁。

\textbf{公式正答：} 5（single）。

\par\textit{原典本文の校正と独自解説は未作成。分類のみの登録項目。}

\section{有限要素法}

\subsection{平成23年度 問3-1}
\label{q:h23-reg-basic-3-1}

\noindent\textbf{原典：} sources/H23/basic.pdf、PDF物理8頁。

\textbf{公式正答：} 2（single）。

\par\textit{原典本文の校正と独自解説は未作成。分類のみの登録項目。}

\subsection{平成24年度 問3-1}
\label{q:h24-reg-basic-3-1}

\noindent\textbf{原典：} sources/H24/basic.pdf、PDF物理8頁。

\textbf{公式正答：} 1（single）。

\par\textit{原典本文の校正と独自解説は未作成。分類のみの登録項目。}

\subsection{平成25年度 問3-3}
\label{q:h25-reg-basic-3-3}

\noindent\textbf{原典：} sources/H25/basic.pdf、PDF物理12頁。

\textbf{公式正答：} 1（single）。

\par\textit{原典本文の校正と独自解説は未作成。分類のみの登録項目。}

\subsection{平成29年度 問3-3}
\label{q:h29-reg-basic-3-3}

\noindent\textbf{原典：} sources/H29/basic.pdf、PDF物理11頁。

\textbf{公式正答：} 4（single）。

\par\textit{原典本文の校正と独自解説は未作成。分類のみの登録項目。}

\subsection{令和2年度 問3-4}
\label{q:r02-reg-basic-3-4}

\noindent\textbf{原典：} sources/R02/basic.pdf、PDF物理12頁。

\textbf{公式正答：} 4（single）。

\par\textit{原典本文の校正と独自解説は未作成。分類のみの登録項目。}

\subsection{令和3年度 問3-3}
\label{q:r03-reg-basic-3-3}

\noindent\textbf{原典：} sources/R03/basic.pdf、PDF物理12頁。

\textbf{公式正答：} 1（single）。

\par\textit{原典本文の校正と独自解説は未作成。分類のみの登録項目。}

\section{材料の応力・変形}

\subsection{平成23年度 問3-2}
\label{q:h23-reg-basic-3-2}

\noindent\textbf{原典：} sources/H23/basic.pdf、PDF物理9頁。

\textbf{公式正答：} 1（single）。

\par\textit{原典本文の校正と独自解説は未作成。分類のみの登録項目。}

\subsection{平成24年度 問3-5}
\label{q:h24-reg-basic-3-5}

\noindent\textbf{原典：} sources/H24/basic.pdf、PDF物理11頁。

\textbf{公式正答：} 4（single）。

\par\textit{原典本文の校正と独自解説は未作成。分類のみの登録項目。}

\subsection{平成25年度 問3-1}
\label{q:h25-reg-basic-3-1}

\noindent\textbf{原典：} sources/H25/basic.pdf、PDF物理10頁。

\textbf{公式正答：} 5（single）。

\par\textit{原典本文の校正と独自解説は未作成。分類のみの登録項目。}

\subsection{平成25年度 問3-5}
\label{q:h25-reg-basic-3-5}

\noindent\textbf{原典：} sources/H25/basic.pdf、PDF物理13頁。

\textbf{公式正答：} 2（single）。

\par\textit{原典本文の校正と独自解説は未作成。分類のみの登録項目。}

\subsection{平成26年度 問3-6}
\label{q:h26-reg-basic-3-6}

\noindent\textbf{原典：} sources/H26/basic.pdf、PDF物理12頁。

\textbf{公式正答：} 2（single）。

\par\textit{原典本文の校正と独自解説は未作成。分類のみの登録項目。}

\subsection{平成27年度 問3-5}
\label{q:h27-reg-basic-3-5}

\noindent\textbf{原典：} sources/H27/basic.pdf、PDF物理10頁。

\textbf{公式正答：} 3（single）。

\par\textit{原典本文の校正と独自解説は未作成。分類のみの登録項目。}

\subsection{平成27年度 問3-6}
\label{q:h27-reg-basic-3-6}

\noindent\textbf{原典：} sources/H27/basic.pdf、PDF物理11頁。

\textbf{公式正答：} 4（single）。

\par\textit{原典本文の校正と独自解説は未作成。分類のみの登録項目。}

\subsection{平成28年度 問3-6}
\label{q:h28-reg-basic-3-6}

\noindent\textbf{原典：} sources/H28/basic.pdf、PDF物理16頁。

\textbf{公式正答：} 4（single）。

\par\textit{原典本文の校正と独自解説は未作成。分類のみの登録項目。}

\subsection{平成29年度 問3-5}
\label{q:h29-reg-basic-3-5}

\noindent\textbf{原典：} sources/H29/basic.pdf、PDF物理13頁。

\textbf{公式正答：} 3（single）。

\par\textit{原典本文の校正と独自解説は未作成。分類のみの登録項目。}

\subsection{平成30年度 問3-6}
\label{q:h30-reg-basic-3-6}

\noindent\textbf{原典：} sources/H30/basic.pdf、PDF物理14頁。

\textbf{公式正答：} 3（single）。

\par\textit{原典本文の校正と独自解説は未作成。分類のみの登録項目。}

\subsection{令和元年度 問3-5}
\label{q:r01-reg-basic-3-5}

\noindent\textbf{原典：} sources/R01/basic.pdf、PDF物理14頁。

\textbf{公式正答：} 4（single）。

\par\textit{原典本文の校正と独自解説は未作成。分類のみの登録項目。}

\subsection{令和元年度再試験 問3-6}
\label{q:r01-ret-basic-3-6}

\noindent\textbf{原典：} sources/R01\_retest/basic.pdf、PDF物理13頁。

\textbf{公式正答：} 2（single）。

\par\textit{原典本文の校正と独自解説は未作成。分類のみの登録項目。}

\subsection{令和2年度 問3-6}
\label{q:r02-reg-basic-3-6}

\noindent\textbf{原典：} sources/R02/basic.pdf、PDF物理14頁。

\textbf{公式正答：} 5（single）。

\par\textit{原典本文の校正と独自解説は未作成。分類のみの登録項目。}

\subsection{令和3年度 問3-4}
\label{q:r03-reg-basic-3-4}

\noindent\textbf{原典：} sources/R03/basic.pdf、PDF物理13頁。

\textbf{公式正答：} 3（single）。

\par\textit{原典本文の校正と独自解説は未作成。分類のみの登録項目。}

\subsection{令和4年度 問3-4}
\label{q:r04-reg-basic-3-4}

\noindent\textbf{原典：} sources/R04/basic.pdf、PDF物理11頁。

\textbf{公式正答：} 2（single）。

\par\textit{原典本文の校正と独自解説は未作成。分類のみの登録項目。}

\subsection{令和5年度 問3-4}
\label{q:r05-reg-basic-3-4}

\noindent\textbf{原典：} sources/R05/basic.pdf、PDF物理10頁。

\textbf{公式正答：} 5（single）。

\par\textit{原典本文の校正と独自解説は未作成。分類のみの登録項目。}

\subsection{令和6年度 問3-4}
\label{q:r06-reg-basic-3-4}

\noindent\textbf{原典：} sources/R06/basic.pdf、PDF物理12頁。

\textbf{公式正答：} 3（single）。

\par\textit{原典本文の校正と独自解説は未作成。分類のみの登録項目。}

\subsection{令和7年度 問3-4}
\label{q:r07-reg-basic-3-4}

\noindent\textbf{原典：} sources/R07/basic.pdf、PDF物理10頁。

\textbf{公式正答：} 2（single）。

\par\textit{原典本文の校正と独自解説は未作成。分類のみの登録項目。}

\section{数値解析・誤差}

\subsection{平成26年度 問3-2}
\label{q:h26-reg-basic-3-2}

\noindent\textbf{原典：} sources/H26/basic.pdf、PDF物理9頁。

\textbf{公式正答：} 3（single）。

\par\textit{原典本文の校正と独自解説は未作成。分類のみの登録項目。}

\subsection{平成27年度 問3-3}
\label{q:h27-reg-basic-3-3}

\noindent\textbf{原典：} sources/H27/basic.pdf、PDF物理9頁。

\textbf{公式正答：} 1（single）。

\par\textit{原典本文の校正と独自解説は未作成。分類のみの登録項目。}

\subsection{令和元年度再試験 問3-3}
\label{q:r01-ret-basic-3-3}

\noindent\textbf{原典：} sources/R01\_retest/basic.pdf、PDF物理11頁。

\textbf{公式正答：} 5（single）。

\par\textit{原典本文の校正と独自解説は未作成。分類のみの登録項目。}

\subsection{令和2年度 問3-3}
\label{q:r02-reg-basic-3-3}

\noindent\textbf{原典：} sources/R02/basic.pdf、PDF物理11頁。

\textbf{公式正答：} 3（single）。

\par\textit{原典本文の校正と独自解説は未作成。分類のみの登録項目。}

\subsection{令和4年度 問3-3}
\label{q:r04-reg-basic-3-3}

\noindent\textbf{原典：} sources/R04/basic.pdf、PDF物理11頁。

\textbf{公式正答：} 4（single）。

\par\textit{原典本文の校正と独自解説は未作成。分類のみの登録項目。}

\subsection{令和5年度 問3-3}
\label{q:r05-reg-basic-3-3}

\noindent\textbf{原典：} sources/R05/basic.pdf、PDF物理10頁。

\textbf{公式正答：} 1（single）。

\par\textit{原典本文の校正と独自解説は未作成。分類のみの登録項目。}

\section{力学系の固有振動}

\subsection{平成23年度 問3-4}
\label{q:h23-reg-basic-3-4}

\noindent\textbf{原典：} sources/H23/basic.pdf、PDF物理11頁。

\textbf{公式正答：} 3（single）。

\par\textit{原典本文の校正と独自解説は未作成。分類のみの登録項目。}

\subsection{平成24年度 問3-2}
\label{q:h24-reg-basic-3-2}

\noindent\textbf{原典：} sources/H24/basic.pdf、PDF物理8頁。

\textbf{公式正答：} 2（single）。

\paragraph{問題の要旨} ばね接続から有効剛性を求め、3系の固有振動数を比較する。

\Needspace{4\baselineskip}\paragraph{確認した条件・記述（独自要約）}\mbox{}\par

\noindent\textbf{spring\_constants：} k,k\par

\noindent\textbf{mass：} m\par

\noindent\textbf{a：} two parallel springs between left wall and mass\par

\noindent\textbf{b：} two springs in series between left wall and mass\par

\noindent\textbf{c：} mass between two fixed walls, one spring on each side\par

\paragraph{問われる核心} 振動・運動

\paragraph{必要知識} ばねの直列・並列合成、微小変位の復元力、固有角振動数

\paragraph{考え方（査読待ち初稿）} aでは復元力が2kxなので有効剛性2k。bでは変形量が加算され1/keff=\allowbreak{}1/k+1/k、よってkeff=\allowbreak{}k/2。cも左右の復元力が同方向に合計され2kx。質量が同じため固有角振動数はsqrt(keff/m)で比較でき、bのみ最小。正答②。

\paragraph{他の選択肢（査読待ち初稿）} ①aと③cはいずれも有効剛性2k。④は最小でないaを、⑤は最小でないcを含む。cの左右の力は打ち消し合わず、変位に対して同じ復元方向に働く。

\paragraph{関連問題} \hyperref[q:h26-reg-basic-3-5]{H26-REG-BASIC-3-5}

\par 難易度 2（校準初稿）、優先度未評価。

\subsection{平成25年度 問3-2}
\label{q:h25-reg-basic-3-2}

\noindent\textbf{原典：} sources/H25/basic.pdf、PDF物理11頁。

\textbf{公式正答：} 3（single）。

\par\textit{原典本文の校正と独自解説は未作成。分類のみの登録項目。}

\subsection{平成26年度 問3-5}
\label{q:h26-reg-basic-3-5}

\noindent\textbf{原典：} sources/H26/basic.pdf、PDF物理11頁。

\textbf{公式正答：} 2（single）。

\par\textit{原典本文の校正と独自解説は未作成。分類のみの登録項目。}

\subsection{平成28年度 問3-5}
\label{q:h28-reg-basic-3-5}

\noindent\textbf{原典：} sources/H28/basic.pdf、PDF物理15頁。

\textbf{公式正答：} 1（single）。

\par\textit{原典本文の校正と独自解説は未作成。分類のみの登録項目。}

\subsection{平成29年度 問3-6}
\label{q:h29-reg-basic-3-6}

\noindent\textbf{原典：} sources/H29/basic.pdf、PDF物理14頁。

\textbf{公式正答：} 5（single）。

\par\textit{原典本文の校正と独自解説は未作成。分類のみの登録項目。}

\subsection{平成30年度 問3-5}
\label{q:h30-reg-basic-3-5}

\noindent\textbf{原典：} sources/H30/basic.pdf、PDF物理14頁。

\textbf{公式正答：} 3（single）。

\par\textit{原典本文の校正と独自解説は未作成。分類のみの登録項目。}

\subsection{令和元年度再試験 問3-5}
\label{q:r01-ret-basic-3-5}

\noindent\textbf{原典：} sources/R01\_retest/basic.pdf、PDF物理12頁。

\textbf{公式正答：} 4（single）。

\par\textit{原典本文の校正と独自解説は未作成。分類のみの登録項目。}

\subsection{令和2年度 問3-5}
\label{q:r02-reg-basic-3-5}

\noindent\textbf{原典：} sources/R02/basic.pdf、PDF物理13頁。

\textbf{公式正答：} 1（single）。

\par\textit{原典本文の校正と独自解説は未作成。分類のみの登録項目。}

\subsection{令和3年度 問3-5}
\label{q:r03-reg-basic-3-5}

\noindent\textbf{原典：} sources/R03/basic.pdf、PDF物理14頁。

\textbf{公式正答：} 4（single）。

\par\textit{原典本文の校正と独自解説は未作成。分類のみの登録項目。}

\subsection{令和4年度 問3-6}
\label{q:r04-reg-basic-3-6}

\noindent\textbf{原典：} sources/R04/basic.pdf、PDF物理13頁。

\textbf{公式正答：} 5（single）。

\par\textit{原典本文の校正と独自解説は未作成。分類のみの登録項目。}

\subsection{令和6年度 問3-5}
\label{q:r06-reg-basic-3-5}

\noindent\textbf{原典：} sources/R06/basic.pdf、PDF物理13頁。

\textbf{公式正答：} 3（single）。

\par\textit{原典本文の校正と独自解説は未作成。分類のみの登録項目。}

\section{未確定}

\subsection{平成23年度 問3-3}
\label{q:h23-reg-basic-3-3}

\noindent\textbf{原典：} sources/H23/basic.pdf、PDF物理10頁。

\textbf{公式正答：} 5（single）。

\par\textit{原典本文の校正と独自解説は未作成。分類のみの登録項目。}

\subsection{平成25年度 問3-4}
\label{q:h25-reg-basic-3-4}

\noindent\textbf{原典：} sources/H25/basic.pdf、PDF物理13頁。

\textbf{公式正答：} 5（single）。

\par\textit{原典本文の校正と独自解説は未作成。分類のみの登録項目。}

\subsection{平成27年度 問3-1}
\label{q:h27-reg-basic-3-1}

\noindent\textbf{原典：} sources/H27/basic.pdf、PDF物理8頁。

\textbf{公式正答：} 2（single）。

\par\textit{原典本文の校正と独自解説は未作成。分類のみの登録項目。}

\subsection{平成27年度 問3-2}
\label{q:h27-reg-basic-3-2}

\noindent\textbf{原典：} sources/H27/basic.pdf、PDF物理8頁。

\textbf{公式正答：} 2（single）。

\par\textit{原典本文の校正と独自解説は未作成。分類のみの登録項目。}

\subsection{平成27年度 問3-4}
\label{q:h27-reg-basic-3-4}

\noindent\textbf{原典：} sources/H27/basic.pdf、PDF物理9頁。

\textbf{公式正答：} 3（single）。

\par\textit{原典本文の校正と独自解説は未作成。分類のみの登録項目。}

\subsection{平成29年度 問3-4}
\label{q:h29-reg-basic-3-4}

\noindent\textbf{原典：} sources/H29/basic.pdf、PDF物理12頁。

\textbf{公式正答：} 2（single）。

\par\textit{原典本文の校正と独自解説は未作成。分類のみの登録項目。}

\subsection{平成30年度 問3-1}
\label{q:h30-reg-basic-3-1}

\noindent\textbf{原典：} sources/H30/basic.pdf、PDF物理11頁。

\textbf{公式正答：} 1（single）。

\par\textit{原典本文の校正と独自解説は未作成。分類のみの登録項目。}

\subsection{平成30年度 問3-3}
\label{q:h30-reg-basic-3-3}

\noindent\textbf{原典：} sources/H30/basic.pdf、PDF物理12頁。

\textbf{公式正答：} 2（single）。

\par\textit{原典本文の校正と独自解説は未作成。分類のみの登録項目。}

\subsection{平成30年度 問3-4}
\label{q:h30-reg-basic-3-4}

\noindent\textbf{原典：} sources/H30/basic.pdf、PDF物理13頁。

\textbf{公式正答：} 2（single）。

\par\textit{原典本文の校正と独自解説は未作成。分類のみの登録項目。}

\subsection{令和元年度 問3-3}
\label{q:r01-reg-basic-3-3}

\noindent\textbf{原典：} sources/R01/basic.pdf、PDF物理13頁。

\textbf{公式正答：} 1（single）。

\par\textit{原典本文の校正と独自解説は未作成。分類のみの登録項目。}

\subsection{令和元年度 問3-6}
\label{q:r01-reg-basic-3-6}

\noindent\textbf{原典：} sources/R01/basic.pdf、PDF物理15頁。

\textbf{公式正答：} 3（single）。

\par\textit{原典本文の校正と独自解説は未作成。分類のみの登録項目。}

\subsection{令和元年度再試験 問3-2}
\label{q:r01-ret-basic-3-2}

\noindent\textbf{原典：} sources/R01\_retest/basic.pdf、PDF物理11頁。

\textbf{公式正答：} 2（single）。

\par\textit{原典本文の校正と独自解説は未作成。分類のみの登録項目。}

\subsection{令和3年度 問3-6}
\label{q:r03-reg-basic-3-6}

\noindent\textbf{原典：} sources/R03/basic.pdf、PDF物理15頁。

\textbf{公式正答：} 5（single）。

\par\textit{原典本文の校正と独自解説は未作成。分類のみの登録項目。}

\subsection{令和4年度 問3-5}
\label{q:r04-reg-basic-3-5}

\noindent\textbf{原典：} sources/R04/basic.pdf、PDF物理12頁。

\textbf{公式正答：} 3（single）。

\par\textit{原典本文の校正と独自解説は未作成。分類のみの登録項目。}

\subsection{令和5年度 問3-1}
\label{q:r05-reg-basic-3-1}

\noindent\textbf{原典：} sources/R05/basic.pdf、PDF物理9頁。

\textbf{公式正答：} 4（single）。

\par\textit{原典本文の校正と独自解説は未作成。分類のみの登録項目。}

\subsection{令和5年度 問3-5}
\label{q:r05-reg-basic-3-5}

\noindent\textbf{原典：} sources/R05/basic.pdf、PDF物理11頁。

\textbf{公式正答：} 2（single）。

\par\textit{原典本文の校正と独自解説は未作成。分類のみの登録項目。}

\subsection{令和5年度 問3-6}
\label{q:r05-reg-basic-3-6}

\noindent\textbf{原典：} sources/R05/basic.pdf、PDF物理12頁。

\textbf{公式正答：} 3（single）。

\par\textit{原典本文の校正と独自解説は未作成。分類のみの登録項目。}

\subsection{令和6年度 問3-3}
\label{q:r06-reg-basic-3-3}

\noindent\textbf{原典：} sources/R06/basic.pdf、PDF物理11頁。

\textbf{公式正答：} 5（single）。

\par\textit{原典本文の校正と独自解説は未作成。分類のみの登録項目。}

\subsection{令和6年度 問3-6}
\label{q:r06-reg-basic-3-6}

\noindent\textbf{原典：} sources/R06/basic.pdf、PDF物理14頁。

\textbf{公式正答：} 3（single）。

\par\textit{原典本文の校正と独自解説は未作成。分類のみの登録項目。}

\subsection{令和7年度 問3-1}
\label{q:r07-reg-basic-3-1}

\noindent\textbf{原典：} sources/R07/basic.pdf、PDF物理8頁。

\textbf{公式正答：} 2（single）。

\par\textit{原典本文の校正と独自解説は未作成。分類のみの登録項目。}

\subsection{令和7年度 問3-5}
\label{q:r07-reg-basic-3-5}

\noindent\textbf{原典：} sources/R07/basic.pdf、PDF物理11頁。

\textbf{公式正答：} 5（single）。

\par\textit{原典本文の校正と独自解説は未作成。分類のみの登録項目。}

\subsection{令和7年度 問3-6}
\label{q:r07-reg-basic-3-6}

\noindent\textbf{原典：} sources/R07/basic.pdf、PDF物理12頁。

\textbf{公式正答：} 1（single）。

\par\textit{原典本文の校正と独自解説は未作成。分類のみの登録項目。}
